%% Per-system sensitivity of the carrier search.  Table fragment generated by
%% persys_r17.py; the caption quotes only macros other generators publish.
%% r30 (write-apx): three sentences above Table F1 and short notes.  The
%% decorrelation caveat is Table 3's; the notes keep only how the column is
%% formed.  Register row 3: the rows of \OmFlagStars are conditional on the
%% adopted position (Appendix A.2).  ★ r30-num: persys_r17.py must print the
%% mark the note names (a section mark, $^{\S}$) on those three rows, or this
%% note must be changed to the mark it prints.
%% r33 (round-32 R1-5): the pilcrow ($\P$, persys_r17 P12, r1_r22 A14) marks the extrapolated-limit
%% systems; the window-to-window dispersion is attached to the interpolated transfers only (rankcharge_r33 NG4).
%% r34 (W3): notes for the crossing-window mark (R2-1, N1) and the
%% decorrelation-corrected EIRP (R2-4, N1).  The mark must be the one
%% persys_r17.py prints (\XwMark, double bar, persys_r17 P13).
%% r35: the redraw paragraph (transfer dispersion over the sample) and the rank-cost paragraph
%% (rankcharge_r33 macros) MOVED HERE from Sec. 4.1, which keeps one pointer sentence for each.  The table
%% notes now say what the pilcrow means (R1-3: extrapolated limit, no dispersion attached) and which
%% window sets AU Mic's printed limit (m7).  NG4 (rankcharge_r33): no sentence here may put
%% \SmcDispLo beside \SmcNEstSysBnd.
%% r36: dec. is now the correction DIVIDED OUT of the adopted limit (R2-8; measured, or imputed = asterisk);
%% the note gives the limits BEFORE it for the >10 % systems and says the correction is a floor.  A new mark
%% (R2-4, persys_r17 must print the mark the note names) flags best windows beyond the largest injected offset.
%% r36 phase 2: the switched-window sentences (\XwSwitchNames, \NtSw*) are GONE (no system switches; N2). The note on
%% \EvAStar's limit block vs its discovery block (N1 hook 2) answers why its limit is deep: the arrays differ, no correction.
\section{Per-system sensitivity}
\label{app:persys}

Table~\ref{tab:persys} gives for each system the limit of its best window
outside the attribution mask (\S\ref{sec:summary}). A bound pair appears
under its primary's name, and the limits at both ends of the calibration
interval are in \texttt{sensitivity/\allowbreak system\_limits.csv} of the
data archive \citep{SIdeposit}. Beside each limit are that window's centre
frequency $\nu$ and channel width $\Delta\nu$, the system's merged coverage
$\Delta\nu_{\rm cov}$, the percentage \emph{dec.} by which atmospheric
decorrelation raises the limit, already divided out of the tabulated value,
and whether blocks of different
sessions cover a common frequency more than a day (d) or a year (yr) apart,
over the census and the \ZcNRecEB{} blocks the recurrence test reads. A
dagger marks the \SmcNMeasSys{} systems whose best window was itself injected
into, \SmcNMeasSysBnd{} of them among the \NSysALimSum{} systems with a bounded
position. The other \SmcNEstSysBnd{} limits of those \NSysALimSum{} are transferred
from the class average. For the \NgNInterpEstSysBnd{} of them whose best window
lies inside the configurations the injections covered, a single limit is
uncertain by the dispersion between injected windows, $\times\SmcDispLo$ to
$\times\SmcDispHi$. A pilcrow ($\P$) marks the \NgNExtrapSys{}
\NgExtrapSysNoun{} whose best window lies outside them, \NgExtrapSys. No
measured dispersion applies to those limits, so we also give the summaries
without them: over the other \NgNSysInterpSum{} systems the median limit is
$\NgSysMedInterp$\,W and \NgNSysEfifteenInterp{} reach $10^{15}$\,W.

The effect of that dispersion on the sample as a whole was found by drawing
each transferred window's factor at random from the measured ones, once per
execution block, over \SvNTrial{} trials. This lowers the median by
\SvRsMedShiftPct{} per cent and leaves \SvRsBlkNLo--\SvRsBlkNHi{} systems at
$10^{15}$\,W in ninety per cent of the redraws. Even applied at its extreme to every
transferred window, the spread could move \SvSprN{} systems across that
level, \SvSprNOut{} out and \SvSprNIn{} in. Each system's probability of
crossing it is given in the data archive.

None of these limits is charged for the spatial rank against the \NCtrl{}
control positions, which orders the crossings for examination and disposes
of none of them (\S\ref{sec:pninetysel}). Charging it would raise the
limits, though not by a single factor. Among the injected windows the cost
runs from $\times\RcWinCostLo$ to $\times\RcWinCostHi$, with a median of
$\times\RcWinCostMed$, and it grows with the brightness of the control ring.
Averaged over the class it is $\times\SensCostLike$, which would move the
median system to $\EirpNinetySysMedARank$\,W. If each window is instead
charged at the cost measured for windows with a similar ring, the median
becomes $\RcSysMedStrat$\,W and \RcNSysEfifteenStrat{} systems still reach
$10^{15}$\,W. In \SensNRingHiUnres{} windows with a very bright ring,
\RcNoRankWins, the rank could not be met at any injected power.
\RcNoRankSysClause.

\begin{table*}
\centering
\caption{\textbf{Per-system sensitivity of the carrier search.} For each of
the \NSysClassA{} systems the fine-channel search covers, the
$\mathrm{EIRP}_{90}$ of its best window, retentions applied, in order of
the limit down the left half and then down the right half. Columns and
marks: Appendix~\ref{app:persys} and the notes below.}
\label{tab:persys}
\footnotesize
%% GENERATED by persys_r17.py -- do not hand-edit.
%% One row per system of the carrier search, two systems per printed line,
%% ordered by the limit.  Every value is computed from the per-window results.
\begin{tabular}{@{}l@{\hspace{4pt}}r@{\hspace{4pt}}r@{\hspace{4pt}}r@{\hspace{4pt}}r@{\hspace{4pt}}r@{\hspace{4pt}}c@{\hspace{10pt}}l@{\hspace{4pt}}r@{\hspace{4pt}}r@{\hspace{4pt}}r@{\hspace{4pt}}r@{\hspace{4pt}}r@{\hspace{4pt}}c@{}}
\hline
System & $\mathrm{EIRP}_{90}$ & $\nu$ & $\Delta\nu$ & $\Delta\nu_{\rm cov}$ & dec. & conf. & System & $\mathrm{EIRP}_{90}$ & $\nu$ & $\Delta\nu$ & $\Delta\nu_{\rm cov}$ & dec. & conf. \\
 & ($10^{14}$\,W) & (GHz) & (kHz) & (GHz) & (\%) &  &  & ($10^{14}$\,W) & (GHz) & (kHz) & (GHz) & (\%) &  \\
\hline
$\epsilon$~Eri & 0.37$^{\diamond}$ & 229.98 & 977 & 2.3 & -- & yr & $\kappa$~CrB & 14$^{\dagger}$ & 230.57 & 977 & 1.7 & 2$^{*}$ & -- \\
LP 876-10 & 0.77 & 345.83 & 488 & 1.7 & 4$^{*}$ & d & HD 170773 & 15 & 230.13 & 488 & 7.4 & 2$^{*}$ & d \\
$\tau$~Cet & 0.79 & 230.55 & 488 & 1.8 & 0.4$^{*}$ & d & HD 61005 & 17 & 230.52 & 244 & 19 & 3$^{*}$ & d \\
HD 69830 & 0.94$^{\dagger}$ & 335.38 & 244 & 3.8 & 0.5 & d & LHS 1140 & 18$^{\S\|}$ & 230.00 & 488 & 1.8 & 0.4$^{*}$ & d \\
AU Mic & 1.1 & 230.56 & 488 & 3.5 & 3$^{*}$ & yr & LP 736-15 & 18 & 345.82 & 977 & 1.7 & 4$^{*}$ & -- \\
61 Vir & 1.2 & 345.79 & 488 & 1.8 & 5$^{*}$ & yr & L 836-122 & 20 & 346.00 & 488 & 1.7 & 1 & -- \\
TRAPPIST-1 & 1.5 & 346.58 & 488 & 1.7 & 5 & d & GJ 2006A & 22 & 346.02 & 488 & 1.7 & 1 & -- \\
AT Mic AB & 1.6 & 346.04 & 488 & 1.7 & 3 & -- & HD 14055 & 23 & 230.11 & 488 & 7.2 & 3 & yr \\
$\beta$~Pic & 1.6$^{\ddagger}$ & 226.85 & 122 & 4.4 & 2$^{*}$ & yr & HD 23484 & 24 & 230.53 & 977 & 1.8 & 0.2 & d \\
HD 38858 & 2.0 & 230.50 & 488 & 1.8 & 3$^{*}$ & yr & CP$-$72 2713 & 25 & 345.12 & 488 & 5.8 & 8 & d \\
GJ 3379 & 2.0$^{\|}$ & 229.96 & 488 & 1.9 & 0.4$^{*}$ & d & TWA 7 & 26$^{\P}$ & 345.99 & 244 & 0.87 & 26$^{*}$ & d \\
HD 207129 & 2.2 & 230.57 & 488 & 1.8 & 3$^{*}$ & yr & CD$-$57 1054 & 26 & 345.99 & 488 & 1.7 & 20 & d \\
HD 139664 & 3.0 & 230.55 & 488 & 1.7 & 3$^{*}$ & -- & UCAC3 176-23654 & 27 & 346.02 & 488 & 1.7 & 4 & -- \\
GJ 14 & 3.4 & 227.22 & 488 & 3.5 & 19 & d & TYC 2211-1309-1 & 27 & 346.00 & 488 & 1.7 & 5 & -- \\
HD 53143 & 3.5 & 230.53 & 977 & 1.8 & 2$^{*}$ & yr & HD 155555C & 28$^{\dagger\S}$ & 346.01 & 488 & 1.7 & 43 & -- \\
HD 202628 & 3.5 & 230.51 & 488 & 1.8 & 4$^{*}$ & yr & PM J03433+1958 & 28$^{\P}$ & 230.13 & 488 & 1.7 & 2$^{*}$ & -- \\
HN Lib & 3.5$^{\dagger}$ & 229.99 & 488 & 1.8 & 0.4$^{*}$ & d & 51 Eri & 28$^{\P}$ & 230.53 & 1953 & 1.7 & 11$^{*}$ & -- \\
HD 10647 & 3.6 & 230.54 & 488 & 6.1 & 4$^{*}$ & yr & Barta 161 12 & 32$^{\S}$ & 346.02 & 488 & 1.7 & 4$^{*}$ & -- \\
$\eta$~Crv & 5.1 & 343.35 & 488 & 7.2 & 27 & yr & UCAC2 20312880 & 33 & 345.96 & 488 & 1.7 & 11 & -- \\
GJ 849 & 6.2 & 230.02 & 488 & 1.9 & 0.4$^{*}$ & d & LP 353-51 & 37 & 346.02 & 488 & 1.7 & 26 & -- \\
TWA 3A & 6.4$^{\dagger}$ & 230.54 & 122 & 0.44 & 35$^{*}$ & d & HD 139084 & 50 & 346.01 & 488 & 1.7 & 3 & -- \\
HD 92945 & 7.7 & 355.11 & 488 & 3.7 & 6$^{*}$ & d & HD 172555 & 57$^{\P}$ & 340.94 & 488 & 5.3 & 200 & -- \\
TX PsA & 8.0 & 346.03 & 488 & 1.7 & 7 & -- & HD 45184 & 59 & 230.53 & 977 & 1.7 & 2 & d \\
HD 285968 & 8.2 & 229.97 & 488 & 1.8 & 0.4$^{*}$ & d & WD 0407+179 & 77$^{\P}$ & 230.13 & 488 & 1.7 & 2$^{*}$ & -- \\
HD 161868 & 9.4 & 227.22 & 488 & 7.4 & 5 & d & HD 48370 & 80$^{\dagger}$ & 230.73 & 122 & 3.3 & 0.4$^{*}$ & yr \\
HD 110411 & 12 & 219.58 & 244 & 4.7 & 2$^{*}$ & yr & HD 31392 & 83 & 230.52 & 977 & 1.7 & 0.9 & d \\
HR 1010 & 12 & 230.53 & 977 & 2.8 & 1 & d & ALMA J153702653$-$33192492 & 93$^{\P}$ & 345.98 & 488 & 14 & 5$^{*}$ & d \\
HIP 490 & 12$^{\dagger}$ & 230.53 & 244 & 0.86 & 3$^{*}$ & -- & HD 89452 & 160 & 230.53 & 977 & 1.7 & 0.7 & d \\
Wolf 219 & 13 & 230.13 & 488 & 1.7 & 2$^{*}$ & d & HD 377 & 290 & 230.52 & 977 & 1.7 & 2 & d \\
HD 107146 & 13 & 355.14 & 488 & 3.7 & 6$^{*}$ & d & HIP 10679 & -- & -- & -- & 0.11 & -- & -- \\
LP 476-207 & 13 & 346.02 & 488 & 1.7 & 4 & -- &  &  &  &  &  &  &  \\
\hline
\end{tabular}


{\footnotesize\raggedright \emph{Notes.} A double dagger marks a row
keeping less than half its channels outside the mask; \NSysALostName, whose
\UmALostWinClause{} inside it, has dashes. A section mark
($\S$) flags \OmFlagStars, whose limits hold only at the adopted position
(Appendix~\ref{app:pbresp}) and which are left out of the median and of
the counts quoted over systems.
\emph{dec.} is the correction divided out of the limit. It is measured
where the delivery records the phase rms, and elsewhere it is imputed as
\DcImputedRule{} and marked with an asterisk. A dash marks the one system with neither,
whose limit is uncorrected. The
correction is itself a floor, because the phase rms is recorded on the
calibrator, so a tabulated limit may still be optimistic by what the
calibrator does not see. For the \DcNBigTab{} systems where decorrelation
exceeds ten per cent\DcBigCondClause, the limits before the correction are,
in units of \DcCorrUnit, \DcNomClause.
A double bar (\XwMark) marks a system whose best-window limit is moved by
the reading at the crossing's own level (\S\ref{sec:pninetysel}). In a window
that holds an unattributed crossing, a carrier reaching only the trigger
would not have stood out, so the limit there is the power at which nine
injected carriers in ten reach the crossing's own statistic. It is read
from the pooled injections, since none of these windows was injected.
The limit of \EvAStar{} comes from its block of \RcSxLimDate, with
\RcSxLimRms\,mJy per \RcSxLimChankHz\,kHz channel, while the crossing that
leads its field lies in a block of the compact array with \RcSxDiscRms\,mJy,
so the two noise levels differ because of the arrays and not because of any
correction.
A pilcrow ($\P$) marks a limit that extrapolates the injection calibration
beyond the configurations that were injected. Such a limit is an
approximate estimate of sensitivity from a model, no injection bounds it,
and no measured dispersion is attached to it (\S\ref{sec:pninetysel}).
\OffMark{} marks a best window in which the star lies beyond the
largest offset at which carriers were injected, \OffInjMaxArcsec$''$, with
a primary-beam response below \OffAcut. Its limit rests on the beam model, and
pointing errors weigh most heavily on it (\S\ref{sec:pninetysel}).\par}
\end{table*}
